% =============================================================================
% fmcw_chirp_beat.tex - SDD Figure 5-1: FMCW principle, multiple targets
% =============================================================================
% Top panel: the transmit chirp (slope S = B/T_C) and three echoes, each the
% same chirp delayed by its round-trip time tau_i = 2 R_i / c. Because every
% echo shares the slope S, the transmit-to-echo gap is a constant beat
% frequency f_bi = S * tau_i.
% Bottom panel: each target therefore appears as one constant beat tone that
% switches on when its echo arrives; a longer range gives a higher beat.
%
% Compile:  pdflatex fmcw_chirp_beat.tex
% Convert:  pdftocairo -svg fmcw_chirp_beat.pdf fmcw_chirp_beat.svg
%           pdftocairo -png -r 200 -singlefile fmcw_chirp_beat.pdf fmcw_chirp_beat
% =============================================================================
\documentclass[border=8pt]{standalone}
%
\input{../style/pmvb-figures.sty}
%
\begin{document}
\begin{tikzpicture}[>=Stealth, font=\sffamily, scale=1.6, transform shape]
  %
  % --- layout constants (letter-only macro names) ---
  \def\Tc{8}        % chirp duration / time-axis length
  \def\Bf{2.8}      % sweep bandwidth / freq-axis height (top panel)
  \def\taua{1.0}    % near-target round-trip delay
  \def\taub{2.0}    % mid-target
  \def\tauc{3.0}    % far-target
  \def\xa{3.5}      % x where beat f_b1 is drawn
  \def\xb{4.7}
  \def\xc{5.9}
  \def\bb{-3.6}     % bottom-panel baseline (y)
  \def\hbA{0.5}     % bottom-panel beat heights (ratio 1:2:3)
  \def\hbB{1.0}
  \def\hbC{1.5}
  \pgfmathsetmacro\Sg{\Bf/\Tc}
  \pgfmathsetmacro\eaY{\Sg*(\Tc-\taua)}
  \pgfmathsetmacro\ebY{\Sg*(\Tc-\taub)}
  \pgfmathsetmacro\ecY{\Sg*(\Tc-\tauc)}
  \pgfmathsetmacro\txa{\Sg*\xa}
  \pgfmathsetmacro\eya{\Sg*(\xa-\taua)}
  \pgfmathsetmacro\txb{\Sg*\xb}
  \pgfmathsetmacro\eyb{\Sg*(\xb-\taub)}
  \pgfmathsetmacro\txc{\Sg*\xc}
  \pgfmathsetmacro\eyc{\Sg*(\xc-\tauc)}
  %
  % =========================== TOP PANEL ====================================
  \draw[pmvbMuted, ->, line width=0.7pt] (-0.1,0) -- (\Tc+0.9,0)
      node[anchor=north east, font=\scriptsize, text=pmvbMuted] {time};
  \draw[pmvbMuted, ->, line width=0.7pt] (0,-0.1) -- (0,\Bf+0.7)
      node[anchor=south west, font=\scriptsize, text=pmvbMuted] {inst.\ frequency};
  \draw[pmvbMuted, line width=0.5pt] (-0.08,\Bf) -- (0.08,\Bf);
  \node[anchor=east, font=\scriptsize, text=pmvbMuted] at (-0.1,\Bf) {$B$};
  %
  % transmit chirp
  \draw[pmvbGreen, line width=1.2pt] (0,0) -- (\Tc,\Bf);
  \node[anchor=south, font=\scriptsize, text=pmvbGreen, rotate=19]
      at (6.7,{\Sg*6.7+0.12}) {TX chirp};
  \node[anchor=south, font=\scriptsize, text=pmvbMuted, rotate=19]
      at (2.1,{\Sg*2.1+0.12}) {$S=B/T_C$};
  %
  % echoes (same slope, delayed)
  \draw[pmvbAmber,  line width=1.0pt, dash pattern=on 4pt off 2.5pt] (\taua,0) -- (\Tc,\eaY);
  \draw[pmvbBlue,   line width=1.0pt, dash pattern=on 4pt off 2.5pt] (\taub,0) -- (\Tc,\ebY);
  \draw[pmvbViolet, line width=1.0pt, dash pattern=on 4pt off 2.5pt] (\tauc,0) -- (\Tc,\ecY);
  \node[anchor=west, font=\scriptsize, text=pmvbMuted] at (\Tc+0.12,\ebY) {echoes};
  %
  % beat-frequency arrows (transmit-to-echo gap), staggered
  \draw[pmvbRed, <->, line width=0.8pt] (\xa,\eya) -- (\xa,\txa);
  \node[anchor=south, font=\tiny, text=pmvbRed] at (\xa,{\txa+0.05}) {$f_{b1}$};
  \draw[pmvbRed, <->, line width=0.8pt] (\xb,\eyb) -- (\xb,\txb);
  \node[anchor=south, font=\tiny, text=pmvbRed] at (\xb,{\txb+0.05}) {$f_{b2}$};
  \draw[pmvbRed, <->, line width=0.8pt] (\xc,\eyc) -- (\xc,\txc);
  \node[anchor=south, font=\tiny, text=pmvbRed] at (\xc,{\txc+0.05}) {$f_{b3}$};
  %
  % tau ticks on the time axis
  \foreach \tk/\lab in {\taua/{$\tau_1$}, \taub/{$\tau_2$}, \tauc/{$\tau_3$}} {
     \draw[pmvbMuted, line width=0.4pt] (\tk,0.07) -- (\tk,-0.07);
     \node[anchor=north, font=\tiny, text=pmvbMuted] at (\tk,-0.05) {\lab};
  }
  \draw[pmvbMuted, line width=0.4pt] (\Tc,0.07) -- (\Tc,-0.07);
  \node[anchor=north, font=\scriptsize, text=pmvbMuted] at (\Tc,-0.07) {$T_C$};
  %
  % =========================== BOTTOM PANEL =================================
  \draw[pmvbMuted, ->, line width=0.7pt] (-0.1,\bb) -- (\Tc+0.9,\bb)
      node[anchor=north east, font=\scriptsize, text=pmvbMuted] {time};
  \draw[pmvbMuted, ->, line width=0.7pt] (0,\bb-0.1) -- (0,{\bb+2.0})
      node[anchor=south west, font=\scriptsize, text=pmvbMuted] {beat freq.};
  %
  \draw[pmvbAmber,  line width=1.0pt] (\taua,\bb) -- (\taua,{\bb+\hbA}) -- (\Tc,{\bb+\hbA});
  \draw[pmvbBlue,   line width=1.0pt] (\taub,\bb) -- (\taub,{\bb+\hbB}) -- (\Tc,{\bb+\hbB});
  \draw[pmvbViolet, line width=1.0pt] (\tauc,\bb) -- (\tauc,{\bb+\hbC}) -- (\Tc,{\bb+\hbC});
  \node[anchor=east, font=\tiny, text=pmvbAmber]  at (-0.1,{\bb+\hbA}) {$f_{b1}$};
  \node[anchor=east, font=\tiny, text=pmvbBlue]   at (-0.1,{\bb+\hbB}) {$f_{b2}$};
  \node[anchor=east, font=\tiny, text=pmvbViolet] at (-0.1,{\bb+\hbC}) {$f_{b3}$};
  \draw[pmvbMuted, line width=0.4pt] (\Tc,{\bb+0.07}) -- (\Tc,{\bb-0.07});
  \node[anchor=north, font=\scriptsize, text=pmvbMuted] at (\Tc,{\bb-0.07}) {$T_C$};
  \node[font=\scriptsize, text=pmvbViolet] at (4.4,{\bb+1.85}) {Beat frequency: one tone per target};
  %
\end{tikzpicture}
\end{document}
